
\begin{tikzpicture}[node distance=0mm and 5mm]
  \node (etape1) [draw, rectangle] {
    \tiny
    \begin{minipage}{0.17\textwidth}
      \inputminted[linenos=false, numbersep=0pt, xleftmargin=0pt, frame=none]
        {ocaml}{exemples/intro/somme_liste_rec_court.ml}
    \end{minipage}
  };
  \node [below=of etape1] {
    \tiny
    \begin{minipage}{0.20\textwidth}
      Fonctions récursives dans un code source au format texte.
    \end{minipage}
  };
  \node (etape2) [draw, rectangle, right=of etape1] {
    \tiny
    \begin{minipage}{0.17\textwidth}
      {\setlength{\fboxsep}{0.8pt}
        \fcolorbox{blue}{white}{\mintinline{ocaml}{let}}
        \fcolorbox{blue}{white}{\mintinline{ocaml}{rec}}
        \fcolorbox{blue}{white}{\mintinline{ocaml}{somme}}
        \fcolorbox{blue}{white}{\mintinline{ocaml}{l}}
        \fcolorbox{blue}{white}{\mintinline{ocaml}{=}}
        \fcolorbox{blue}{white}{\mintinline{ocaml}{match}}
        \fcolorbox{blue}{white}{\mintinline{ocaml}{l}}
        \fcolorbox{blue}{white}{\mintinline{ocaml}{with}}
        \fcolorbox{blue}{white}{\mintinline{ocaml}{|}}
        \fcolorbox{blue}{white}{\mintinline{ocaml}{[}}
        \fcolorbox{blue}{white}{\mintinline{ocaml}{]}}
        \fcolorbox{blue}{white}{\mintinline{ocaml}{->}}
        \fcolorbox{blue}{white}{\mintinline{ocaml}{0}}
        \fcolorbox{blue}{white}{\mintinline{ocaml}{|}}
        \fcolorbox{blue}{white}{\mintinline{ocaml}{x}}
        \fcolorbox{blue}{white}{\mintinline{ocaml}{::}}
        \fcolorbox{blue}{white}{\mintinline{ocaml}{q}}
        \fcolorbox{blue}{white}{\mintinline{ocaml}{->}}
        \fcolorbox{blue}{white}{\mintinline{ocaml}{x}}
        \fcolorbox{blue}{white}{\mintinline{ocaml}{+}}
        \fcolorbox{blue}{white}{\mintinline{ocaml}{somme}}
        \fcolorbox{blue}{white}{\mintinline{ocaml}{q}}
      }
    \end{minipage}
  };
  \node [below=of etape2] {
    \tiny
    Lexèmes
  };
  \node (etape3) [draw, rectangle, right=of etape2] {
      \tiny
      TODO
      %\begin{minipage}{0.5\textwidth}
      %  \begin{tikzpicture}%[
      %      %level 1/.style={sibling distance=5mm},
      %      %level 2/.style={sibling distance=5mm},
      %    %]
      %    \node [draw, rectangle] {\mintinline{ocaml}{let rec}}
      %      child { node [draw, rectangle] {\mintinline{ocaml}{somme}} }
      %      child { node [draw, rectangle] {\mintinline{ocaml}{l}} }
      %      child { node [draw, rectangle] {\mintinline{ocaml}{match}}
      %        child { node [draw, rectangle] {l} }
      %          child { node [draw, rectangle] {[]}
      %          child { node [draw, rectangle] {0}}
      %        }
      %        child { node [draw, rectangle] {+}
      %          child { node [draw, rectangle] {x} }
      %          child { node [draw, rectangle] {somme}
      %            child { node [draw, rectangle] {q} }
      %          }
      %        }
      %      };
      %  \end{tikzpicture}
      %\end{minipage}
  };
  \node [below=of etape3] {
    \tiny
    Arbre
  };
  \node (etape4) [draw, rectangle, right=of etape3] {
      \tiny
      TODO
  };
  \node [below=of etape4] {
    \tiny
    Arbre de la boucle
  };
  \node (etape5) [draw, rectangle, right=of etape4] {
    \tiny
    \begin{minipage}{0.24\textwidth}
      \inputminted[linenos=false, numbersep=0pt, xleftmargin=0pt, frame=none]
        {ocaml}{exemples/intro/somme_liste_boucle_court.ml}
    \end{minipage}
  };
  \node [below=of etape5] {
    \tiny
    Code source de la boucle
  };

  \draw (etape1) edge[->] node[right] {} (etape2);
  \draw (etape2) edge[->] node[right] {} (etape3);
  \draw (etape3) edge[->] node[right] {} (etape4);
  \draw (etape4) edge[->] node[right] {} (etape5);
\end{tikzpicture}
